\documentclass{report}
\usepackage{amsmath,amssymb}
\setlength{\parindent}{0mm}
\setlength{\parskip}{1em}
\begin{document}
\begin{center}
\rule{6in}{1pt} \
{\large Martin Takac \\
{\bf Efficiency of Randomized Coordinate Descent Methods on Minimization Problems with a Composite Objective Function }}

2 Cameron Terrace \\ Edunburgh \\ EH16 5LD \\ Scotland
\\
{\tt martin.taki@gmail.com}\\
Peter Richtarik\\
Martin Takac\end{center}

In this work we develop a randomized block-coordinate descent method for
minimizing the sum of a smooth and a simple nonsmooth block-separable
convex function and establish iteration complexity bounds. This extends
recent results of Nesterov (Efficiency of coordinate descent methods on
huge-scale optimization problems 2010), which cover the smooth case, to
composite minimization, while improving the complexity and simplifying
the analysis. In the smooth case we allow for arbitrary norms and
probability vectors.

Using both synthetic and real data, we demonstrate numerically that the
method is able to solve various optimization problems with a billion
variables. Such problems can be found, for example, in Compressed Sensing
(lasso), Statistics (group lasso), Machine Learning (L1-regularized
logistic or L2 loss function) and Engineering (truss topology design).

For the L1-regularized least squares problem we implement a
GPU-accelerated parallel version of our algorithm (CUDA) and observe
speedups of up to two orders of magnitude when compared with an efficient
serial code (C).


\end{document}
